\documentclass[12pt]{article}
\usepackage{amsmath}
\usepackage{amssymb}
\usepackage{amsthm}
\usepackage{geometry}
\geometry{margin=1.4in}
\usepackage{hyperref}
\usepackage{setspace}
\setstretch{1.4}

\newtheorem{theorem}{Theorem}
\newtheorem{conjecture}[theorem]{Conjecture}
\theoremstyle{definition}
\newtheorem{definition}[theorem]{Definition}
\theoremstyle{remark}
\newtheorem{remark}[theorem]{Remark}

\title{The Field with One Element}
\author{Diego Rinc\'on\\
\small Independent}
\date{}

\begin{document}

\maketitle

\begin{abstract}
$\mathbb{F}_1$ does not exist.
But it should.
If it did, $\mathrm{Spec}\,\mathbb{Z}$ would be a curve over it,
the Frobenius would exist globally,
and the Riemann Hypothesis would follow from the Weil conjectures.
The field with one element is the name of the missing grammar.
\end{abstract}

\section{The Idea}

A field must have at least two elements: $0$ and $1$.
The field axioms require additive and multiplicative inverses,
and $0 \neq 1$.
So $|\mathbb{F}| \geq 2$.
There is no field with one element.

And yet.

Consider the analogy between $\mathbb{Z}$ and $\mathbb{F}_q[t]$.
Both are principal ideal domains.
Both have a theory of primes (integers vs irreducible polynomials).
Both have zeta functions:
$$\zeta_{\mathbb{Z}}(s) = \prod_p \frac{1}{1-p^{-s}},\qquad
  \zeta_{\mathbb{F}_q[t]}(s) = \prod_{f\ \mathrm{irred}} \frac{1}{1-|f|^{-s}}.$$

Over $\mathbb{F}_q$, the ring $\mathbb{F}_q[t]$ is the coordinate ring
of the affine line $\mathbb{A}^1_{\mathbb{F}_q}$.
The integers $\mathbb{Z}$ should be the coordinate ring of
$\mathbb{A}^1_{\mathbb{F}_1}$: the affine line over the field with one element.

If $\mathbb{F}_1$ existed, $\mathrm{Spec}\,\mathbb{Z}$ would be a curve.
Not just a scheme. A curve.
With a smooth compactification: $\overline{\mathrm{Spec}\,\mathbb{Z}}$.
And over a curve, the Weil machinery applies.

\section{What $\mathbb{F}_1$ Would Give}

Over a smooth projective curve $C/\mathbb{F}_q$,
the Riemann Hypothesis holds: the eigenvalues of Frobenius
on $H^1_{\rm\acute{e}t}(C, \mathbb{Q}_\ell)$ have absolute value $q^{1/2}$.
This is proved.

The Riemann Hypothesis for $\zeta(s)$
is the Riemann Hypothesis for the curve $\overline{\mathrm{Spec}\,\mathbb{Z}}$
over $\mathbb{F}_1$.
The zeros of $\zeta(s)$ are the eigenvalues of the $\mathbb{F}_1$-Frobenius
on $H^1_{\mathbb{F}_1}(\overline{\mathrm{Spec}\,\mathbb{Z}}, \mathbb{Q})$.
By the Weil bound: absolute value $1^{1/2} = 1$,
corresponding to $\mathrm{Re}(s) = 1/2$.

This is not a proof.
It is a translation.
The proof would require:
\begin{enumerate}
\item A rigorous definition of $\mathbb{F}_1$ and geometry over it.
\item A construction of $H^1_{\mathbb{F}_1}$.
\item A proof that the Weil bound holds in this setting.
\end{enumerate}
None of the three is complete.
All three are active.

\section{Approaches}

\textbf{Smirnov (1992).}
Noted the analogy between $\mathbb{Z}$ and $\mathbb{F}_q[t]$
and proposed that $\mathbb{F}_1$ should be the ``field of characteristic 1.''
Did not define it.

\textbf{Manin (1995).}
Proposed that $\mathrm{Spec}\,\mathbb{F}_1 = \{*\}$ is the absolute point,
and that all algebraic geometry is relative to it.
Introduced the slogan: finite sets are vector spaces over $\mathbb{F}_1$.

\textbf{Connes--Consani (2010--).}
Defined geometry over $\mathbb{F}_1$ using $\Lambda$-rings and semirings.
Connected $\mathbb{F}_1$ to the Bost--Connes system \cite{bostconnes}:
the symmetry group $\hat{\mathbb{Z}}^* \cong \mathrm{Gal}(\mathbb{F}_1^{\rm ab}/\mathbb{F}_1)$
acts on the KMS states at $\beta > 1$.
The base change from $\mathbb{F}_1$ to $\mathbb{Z}$ is the functor
$\mathbb{Z} \otimes_{\mathbb{F}_1} (-)$.

\textbf{To\"{e}n--Vaqu\'{e} (2009).}
Defined $\mathbb{F}_1$ via monoids: an $\mathbb{F}_1$-algebra is a pointed monoid.
A scheme over $\mathbb{F}_1$ is a sheaf of monoids.
$\mathbb{Z} = \mathbb{F}_1[x, x^{-1}, \ldots]$ in this language.

\textbf{Borger (2009).}
$\Lambda$-rings: $\mathbb{F}_1$-algebras are rings with a compatible
family of Frobenius lifts $\psi^p$ for all primes $p$.
The Adams operations give the $\mathbb{F}_1$-structure.
$\mathrm{Spec}\,\mathbb{Z}$ carries a natural $\Lambda$-structure.
This is the most developed approach.

\section{The Curve}

If $\mathrm{Spec}\,\mathbb{Z}$ is a curve over $\mathbb{F}_1$,
its compactification must include the archimedean place.

The archimedean place is the point at infinity.
This is Arakelov geometry \cite{arakelov}:
compactifying $\mathrm{Spec}\,\mathbb{Z}$
by adjoining the archimedean place.

The curve $\overline{\mathrm{Spec}\,\mathbb{Z}}$ has:
\begin{itemize}
\item One point for each prime $p$: the closed point $p$.
\item One point at infinity: the archimedean place.
\item Genus: unknown. Conjectured to be $0$ (like $\mathbb{P}^1$).
\end{itemize}

If the genus is $0$, the Riemann Hypothesis would say the zeros of $\zeta$
have the right absolute value for a genus-$0$ curve,
consistent with the known analytic results.

\section{What the $\mathbb{F}_1$ Grammar Would Be}

In the language of the grammar paper \cite{grammar}:

The current grammar of arithmetic is the cohomological grammar
of schemes over $\mathbb{Z}$.
This grammar has local Frobenii ($\varphi_p$ at each prime)
but no global Frobenius.
It is locally complete and globally incomplete.

The $\mathbb{F}_1$ grammar would be the cohomological grammar
of schemes over $\mathbb{F}_1$.
In this grammar, $\mathrm{Spec}\,\mathbb{Z}$ is a curve,
and the global Frobenius is the $\mathbb{F}_1$-Frobenius:
a single canonical symmetry acting on the whole curve.

The local-to-global passage \cite{one_wall}
would be solved by the new base:
local Frobenii at each prime cohere globally
because they are all pullbacks of the $\mathbb{F}_1$-Frobenius.

The wall would not be climbed.
The base would be changed.
From the new base, the wall does not appear.

\begin{remark}[What does not yet exist]
The $\mathbb{F}_1$ grammar is not yet a grammar.
It is a sketch of one.
The definition of $\mathbb{F}_1$ is not agreed upon.
The cohomology theory $H^1_{\mathbb{F}_1}$ has not been constructed.
The Weil bound has not been proved in any $\mathbb{F}_1$ setting.
This paper describes what would be true if the grammar existed.
It is not a proof.
It is an anatomy of the absence.
\end{remark}

\begin{thebibliography}{9}

\bibitem{bostconnes}
Rinc\'on, D. (2026).
Bost--Connes: Zeta as Partition Function. Preprint.

\bibitem{grammar}
Rinc\'on, D. (2026).
Mathematics as Cohomological Grammar. Preprint.

\bibitem{one_wall}
Rinc\'on, D. (2026).
One Wall: The Local-to-Global Conjecture. Preprint.

\bibitem{frobenius}
Rinc\'on, D. (2026).
The Missing Frobenius. Preprint.

\bibitem{arakelov}
Rinc\'on, D. (2026).
The Weil Conjectures: Proof of Concept for RH. Preprint.

\bibitem{connes}
Rinc\'on, D. (2026).
Noncommutative Geometry and the Riemann Hypothesis. Preprint.

\end{thebibliography}

\end{document}
